\documentclass[a4paper]{article}

\usepackage[margin=1in]{geometry}
\usepackage{amsmath,amsfonts,mathtools,amsthm,amssymb}
\usepackage[l2tabu,orthodox]{nag}
\usepackage{microtype} % fixes spacing or whatever
\usepackage{enumitem}
\usepackage{tikz-cd}
\usepackage{xcolor}
\usepackage{physics}
\usepackage{mathrsfs}
\usepackage{cancel}
% \usepackage{libertine,libertinust1math}
\usepackage{eucal}
\usepackage[toc,page]{appendix}
\usepackage{pgfplots}
\pgfplotsset{compat=1.18}

\newtheorem{proposition}{Proposition}
\newtheorem{theorem}{Theorem}
\newtheorem{lemma}{Lemma}
\newtheorem{corollary}{Corollary}

\theoremstyle{definition}
\newtheorem{remark}{Remark}
\newtheorem{definition}{Definition}
\newtheorem{example}{Example}

% \usepackage[sc,osf]{mathpazo}   % With old-style figures and real smallcaps.
% \linespread{1.025}              % Palatino leads a little more leading
% % Euler for math and numbers
% \usepackage[euler-digits,small]{eulervm}
% \AtBeginDocument{\renewcommand{\hbar}{\hslash}}

\usepackage{etoolbox}
\usepackage{dynkin-diagrams}
\def\row#1/#2!{\dynkin{#1}{#2}\\}
\newcommand{\tble}[1]{
   \renewcommand*\do[1]{\row##1!}
   \[
      \begin{array}{ll}\docsvlist{#1}\end{array}
   \]
}

\allowdisplaybreaks

\renewcommand{\epsilon}{\varepsilon}
% \renewcommand{\phi}{\varphi}

\DeclareMathAlphabet{\mathpzc}{OT1}{pzc}{m}{it}

\newcommand{\goth}[1]{\mathfrak{#1}}
\newcommand{\red}[1]{#1_{\text{red}}}
\newcommand{\ad}[1]{#1_{\text{Ad}}}
\newcommand{\reg}[1]{#1_{\text{reg}}}
\newcommand{\sh}[1]{#1^{\text{sh}}}
\newcommand{\cpdv}[2]{\cfrac{\partial #1}{\partial #2}}

\newcommand{\fA}{\mathcal{A}}
\newcommand{\fB}{\mathcal{B}}
\newcommand{\fC}{\mathcal{C}}
\newcommand{\fD}{\mathcal{D}}
\newcommand{\fE}{\mathcal{E}}
\newcommand{\fF}{\mathcal{F}}
\newcommand{\fG}{\mathcal{G}}
\newcommand{\fH}{\mathcal{H}}
\newcommand{\fI}{\mathcal{I}}
\newcommand{\fJ}{\mathcal{J}}
\newcommand{\fK}{\mathcal{K}}
\newcommand{\fL}{\mathcal{L}}
\newcommand{\fM}{\mathcal{M}}
\newcommand{\fN}{\mathcal{N}}
\newcommand{\fO}{\mathcal{O}}
\newcommand{\fP}{\mathcal{P}}
\newcommand{\fQ}{\mathcal{Q}}
\newcommand{\fR}{\mathcal{R}}
\newcommand{\fS}{\mathcal{S}}
\newcommand{\fT}{\mathcal{T}}
\newcommand{\fX}{\mathcal{X}}
\newcommand{\fY}{\mathcal{Y}}
\newcommand{\fZ}{\mathcal{Z}}

\newcommand{\be}{\mathbf{e}}
\newcommand{\bx}{\mathbf{x}}
\newcommand{\bone}{\mathbf{1}}
\newcommand{\cx}{\bar{x}}
\newcommand{\cy}{\bar{y}}
\newcommand{\vx}{\vec{x}}
\newcommand{\vsigma}{\vec{\sigma}}
\newcommand{\tzeta}{\tilde{\zeta}}

\newcommand{\gm}{\goth{m}}
\newcommand{\gp}{\goth{p}}
\newcommand{\gF}{\goth{F}}
\newcommand{\gU}{\goth{U}}
\newcommand{\gV}{\goth{V}}

\newcommand{\A}{\mathbf{A}}
\newcommand{\C}{\mathbf{C}}
\newcommand{\F}{\mathbf{F}}
\newcommand{\G}{\mathbf{G}}
\newcommand{\bH}{\mathbf{H}}
\newcommand{\N}{\mathbf{N}}
\newcommand{\PP}{\mathbf{P}}
\newcommand{\Q}{\mathbf{Q}}
\newcommand{\R}{\mathbf{R}}
\newcommand{\Z}{\mathbf{Z}}

\newcommand{\dlog}{\dd\log}

\newcommand\srestr[2]{{\left.\kern-\nulldelimiterspace #1\vphantom{\small|} \right|_{#2}}}
\newcommand\restr[2]{{\left.\kern-\nulldelimiterspace #1 \vphantom{\big|} \right|_{#2}}}
\newcommand\gsec[2]{\Gamma({#1}, {#2})}
\newcommand\gsecx[1]{\Gamma(X, {#1})}

\DeclareMathOperator{\chr}{char}
\DeclareMathOperator{\rProj}{\mathbf{Proj}}
\DeclareMathOperator{\rSpec}{\mathbf{Spec}}
\DeclareMathOperator{\rHom}{\mathbf{Hom}}
\DeclareMathOperator{\rExt}{\mathbf{Ext}}
\DeclareMathOperator{\id}{id}
\DeclareMathOperator{\Frac}{Frac}
\DeclareMathOperator{\rk}{rank}
\DeclareMathOperator{\deter}{det}
\DeclareMathOperator{\pic}{Pic}
\DeclareMathOperator{\cacl}{CaCl}
\DeclareMathOperator{\trd}{tr.d.}
\DeclareMathOperator{\cl}{Cl}
\DeclareMathOperator{\depth}{depth}
\DeclareMathOperator{\codim}{codim}
\DeclareMathOperator{\Div}{Div}
\DeclareMathOperator{\coker}{coker}
\DeclareMathOperator{\len}{length}
\DeclareMathOperator{\height}{height}
\DeclareMathOperator{\supp}{Supp}
\DeclareMathOperator{\proj}{Proj}
\DeclareMathOperator{\im}{im}
\DeclareMathOperator{\Hom}{Hom}
\DeclareMathOperator{\Der}{Der}
\DeclareMathOperator{\spec}{Spec}
\DeclareMathOperator{\Aut}{Aut}
\DeclareMathOperator{\ch}{char}
\DeclareMathOperator{\tor}{Tor}
\DeclareMathOperator{\Ann}{Ann}
\DeclareMathOperator{\Syl}{Syl}
\DeclareMathOperator{\Sym}{Sym}
\DeclareMathOperator{\GL}{GL}
\DeclareMathOperator{\SL}{SL}
\DeclareMathOperator{\Stab}{Stab}
\DeclareMathOperator{\Perm}{Perm}
\DeclareMathOperator{\Orb}{Orb}
\DeclareMathOperator{\Gal}{Gal}
\DeclareMathOperator{\Supp}{Supp}
\DeclareMathOperator{\Ext}{Ext}
\DeclareMathOperator{\Tor}{Tor}
\DeclareMathOperator{\PGL}{PGL}
\DeclareMathOperator{\hd}{hd}
\DeclareMathOperator{\pd}{pd}
\DeclareMathOperator{\SO}{SO}
\DeclareMathOperator{\SU}{SU}
\DeclareMathOperator{\Sp}{Sp}
\DeclareMathOperator{\Oth}{O}
\DeclareMathOperator{\Uni}{U}
\DeclareMathOperator{\evf}{ev}
\DeclareMathOperator{\cH}{H}
\DeclareMathOperator{\dR}{R}
\DeclareMathOperator{\End}{End}
\DeclareMathOperator{\Hasse}{Hasse}
\DeclareMathOperator{\Num}{Num}

\author{James Lee}

% Solarized
% \pagecolor[RGB]{0,20,26}
% \color[RGB]{191,191,191}

% Sephia
% \pagecolor[RGB]{249,239,220}

% Grey
% \pagecolor[RGB]{200, 200, 200}

% Black
% \pagecolor[RGB]{0,0,0}
% \color[RGB]{255,255,255}


\title{Chapter 3, Section 9}

\begin{document}

\maketitle

\begin{enumerate} [label=\textbf{\arabic*.}, leftmargin=0em]

  \item[\textbf{3.}] Some examples of flatness and nonflatness.
        \begin{itemize}
          \item[(a)] If $f : X \to Y$ is a finite surjective morphism of nonsingular varieties over an algebraically closed field $k$, then $f$ is flat.
          \item[(b)] Let $X$ be a union of two planes meeting at a point, each of which maps isomorphically to a plane $Y$. Show that $f$ is not flat. For example, let $Y = \spec{k[x, y]}$ and $$X = \spec{k[x, y, z, w]} / (z, w) \cap (x + z, y + w).$$
          \item[(c)] Again let $Y = \spec{k[x, y]}$, but take $X = \spec{k[x, y, z, w] / (z^2, zw, w^2, xz - yw)}$. Show that $X_\text{red} \cong Y$, $X$ has no embedded points, but that $f$ is not flat.
        \end{itemize}

        \begin{proof} $ $ \vspace{0pt}
          \begin{itemize} [leftmargin=0cm]
            \item[(a)] By Lemma 1, we reduce to the case $X = \spec{B}$ and $Y = \spec{A}$, with $A$ a regular local ring, and $B$ a finite $A$-algebra containing $A$ such that $B_\goth{n}$ is a regular local ring for all maximal ideals $\goth{n} \in \spec{B}$.
                  Then $f$ is flat if and only if $B$ is a projective $A$-module.
                  We use the criterion of (6.12A); that is, we want to show $\depth{B} = \dim{A}$.
                  Let $\goth{m}$ be the maximal ideal of $A$, and let $n = \dim{A}$.
                  It suffices to show $\depth{B} \geq n$ by (6.11A), i.e., we show there exists an $A$-regular sequence for $B$ of length $n$.

                  Regular local rings are integrally closed, so $A$ is integrally closed, and $B$ is integral over $A$, so $\dim{B} = \dim{A}$ by (A.M., 11.26).
                  This implies the following: pick any maximal ideal $\goth{n}$ of $B$, then $B_\goth{n}$ is a regular local ring and $\dim{B_\goth{n}} = \dim{B}$ (A.M., 11.27); in particular, it is a Cohen-Macaulay ring, so a sequence of $n$ elements $x_1, \dots, x_n$ in $B_\goth{n}$ is a $B_\goth{n}$-regular sequence for $B_\goth{n}$ if and only if $B_{\goth{n}}/(x_1, \dots, x_n) = \dim{B} - n = 0$ (II, 8.21A).
                  By the inclusions $A \subseteq B \subseteq B_\goth{n}$, a sequence of elements in $A$ is an $A$-regular sequence for $B$ if and only if it is a $B_\goth{n}$-regular sequence for $B_\goth{n}$.
                  Thus, we further reduce to showing there exists a $B_\goth{n}$-regular sequence contained in $A$ for $B_\goth{n}$ of length $n$.

                  Since $A$ is a regular local ring, there exist $n$ elements $x_1, \dots, x_n \in A$ that form a minimal set of generators of $\goth{m}$ and an $A$-regular sequence for $A$.
                  Viewing $x_1, \dots, x_n$ as elements in $B_\goth{n}$, the elements $x_1, \dots, x_n$ generate the extension of $\goth{m}$ in $B$, so $B_\goth{n} / (x_1, \dots, x_n) \cong B_\goth{n} / \goth{m}B_\goth{n}$, which is an Artin local ring (A.M., 5.8) and has Krull dimension equal to zero.
                  Thus, $x_1, \dots, x_n$ is the desired $B_\goth{n}$-regular sequence for $B_\goth{n}$.

            \item[(b)] We establish the following notation
                  \begin{align*}
                    I & = (z, w) \cap (x + z, y + w) \subseteq k[x, y, z, w], \\
                    Q & = (x, y) \subseteq k[x, y].
                  \end{align*}
                  Let $K = k(x, y)$ be the function field of $Y$.
                  Following the proof of (a), we show
                  \begin{equation}
                    \dim_k k[x, y, z, w]/(I + Q) \neq \dim_K K[z, w]/I.
                  \end{equation}
                  We start with the left hand side of (1) by computing the image of $I + Q$ under the quotient $k[x, y, z, w] \to k[z, w]$ with kernel $Q$.
                  It is not hard to see
                  \begin{equation*}
                    (xz + z^2, yz + zw, xw + zw, yw + w^2) \subseteq I \subseteq (x, y, z, w)^2, \quad (z, w)^2 + Q \subseteq I + Q \subseteq (x, y, z, w)^2 + Q,
                  \end{equation*}
                  and the image of $(x, y, z, w)^2 + Q$ in $k[z, w]$ is $(z, w)^2$, so the image of $I + Q$ is also $(z, w)^2$, i.e.,
                  \begin{equation*}
                    k[x, y, z, w]/(I + Q) \cong k[z, w]/(z, w)^2.
                  \end{equation*}
                  Hence, we have
                  \begin{equation*}
                    \dim_k k[x, y, z, w]/(I + Q) = 3.
                  \end{equation*}
                  It remains to compute the right hand side of (1).
                  The $k[x, y]$-module $k[x, y, z, w]/I$ can be generated by $1, z, w$, so it spans $K[z, w]/I$ over $K$, and $wx - zy$ is in $I$, which implies $z$ and $w$ are $K$-linearly dependent in $K[z, w]/I$.
                  Hence, the $K$-dimension of $K[z, w]/I$ is at most 2.

            \item[(c)] Let $I = (z^2, zw, w^2, xz - yw)$.
                  It suffices to show $$\sqrt{I} = (z, w).$$
                  The generators of $I$ are also in $(z, w)$, a prime ideal, so we have inclusion in one direction, and the left-hand side contains $z, w$, so we have inclusion in the other direction.
                  Also, $I$ is a $(z, w)$-primary ideal, so it can only have one associated point (A.M., 4.5).
                  Hence, $X$ has no embedded points.
                  Lastly, by the same reasoning in (b), we have
                  \begin{equation*}
                    \dim_k k[x, y, z, w]/(I + Q) = 3, \quad \dim_K K[z, w]/I \leq 2.
                  \end{equation*}
          \end{itemize}
        \end{proof}

  \item[\textbf{4.}] \textit{Open Nature of Flatness.} Let $f : X \to Y$ be a morphism of finite type of Noetherian schemes. Then $\{ x \in X \mid \text{$f$ is flat at $x$} \}$ is an open subset of $X$ (possibly empty).

        \begin{proof}
          The question is local on $X$, so we can reduce to the affine case.
          \begin{lemma}
            Let $A$ be a Noetherian ring, $B$ an $A$-algebra of finite type, and $\goth{q}$ a prime ideal of $B$.
            If $B_\goth{q}$ is flat over $A$, then there exists $g \in B - \goth{q}$ such that $B_g$ is flat over $A$.
          \end{lemma}
          We remark that by Lemma 1, $B_\goth{q}$ is flat over $A$ if and only if $B_\goth{q}$ is flat over $A_\goth{p}$, where $\goth{p}$ is the contraction of $\goth{q}$ in $A$.
          Let $E$ be the flat locus of $f$.
          Since $X = \spec{B}$ is a Noetherian space, we can use the following criterion:
          \begin{lemma} [A.M., Ex. 7.22]
            Let $X$ be a Noetherian space, and let $E$ be a subset of $X$.
            Then $E$ is open in $X$ if and only if, for each irreducible closed subset $X_0 \subseteq X$, either $E \cap X_0 = \emptyset$, or else $E \cap X_0$ contains a nonempty open subset of $X_0$.
          \end{lemma}
          In our case, $E$ is open in $X = \spec{B}$ if and only if, for each prime ideal $\goth{q}$ of $B$, either $E \cap V(\goth{q}) = \emptyset$, or else there exists $g \in B - \goth{q}$ such that $V(\goth{q}) \cap X_g \subseteq E \cap V(\goth{q})$.
          Notice that if $E \cap V(\goth{q})$ is nonempty, then $\goth{q} \in E$, because if there exists a prime ideal $\goth{q}' \supset \goth{q}$ such that $B_{\goth{q}'}$ is flat over $A$, then $B_\goth{q}$ is a localization of $B_{\goth{q}'}$, so $B_{\goth{q}}$ is also flat over $A$ by transitivity of flatness.
          Thus, we can replace the condition $V(\goth{q}) \cap E$ is nonempty with $\goth{q}$ is in $E$.

          To proceed, we need a couple preliminary algebraic results:

          \begin{lemma} [Matsumura, 22.3]
            Let $A$ be a Noetherian ring, $I$ an ideal of $A$, and $M$ an $A$-module.
            If $\bigcap_{n = 1}^\infty I^n(\goth{a} \otimes_A M) = 0$ for any ideal $\goth{a}$ of $A$, then the following are equivalent:
            \begin{itemize}
              \item[(i)] $M$ is flat over $A$;
              \item[(ii)] $M/IM$ is flat over $A / I$, and $\tor_1^A(M, A / I) = 0$.
              \item[(iii)] $M / I^nM$ is flat over $A / I^n$ for every $n \geq 1$.
            \end{itemize}
          \end{lemma}
          The main application of Lemma 5 is to the following case: suppose $B$ is a Noetherian $A$-algebra, and $I$ is an ideal in $A$ such that $IB$ is contained in the Jacobson radical of $B$.
          Recall that the radical of a ring is the intersection of all maximal ideals.
          If $\goth{a}$ is any ideal in $A$, then $\goth{a} \otimes_A B$ is a finitely generated $B$ module, so (A.M., 10.19) gives us
          \begin{equation*}
            \bigcap_{n = 1}^\infty I^n(\goth{a} \otimes_A B) = \bigcap_{n = 1}^\infty (IB)^n (\goth{a} \otimes_A B) = 0.
          \end{equation*}
          Hence, $B$ satisfies the hypothesis of Lemma 5.

          We proceed with a partial proof of Lemma 5.
          The implications $(i) \implies (ii) \implies (iii)$ are clear and holds without any special assumptions on $A, I$, or $M$, so we just prove $(iii) \implies (i)$.
          See (Matsumura, 22.3, 24.3) for more details.

          Denote $A_n = A / I^n$ and $M_n = M / I^n M$, and fix an ideal $\goth{a}$ of $A$.
          We want to show $\goth{a} \otimes_A M \xrightarrow{j} M$ has zero kernel.
          By the assumption $\bigcap_n I^n(\goth{a} \otimes_A M) = 0$, it suffices to show $\ker{j} \subseteq I^n(\goth{a} \otimes_A M)$ for all $n$.
          By Krull's Theorem (A.M., 10.11) or (3.1A), for ecah $n$ there exists a sufficiently large integer $k > n$ such that $I^k \cap \goth{a} \subseteq I^n \goth{a}$.
          Consider the commutative diagram
          \[ \begin{tikzcd}
              \goth{a} \otimes_A M \arrow[r, "f"] \arrow[d, "j"'] & (\goth{a} / I^k \cap \goth{a}) \otimes_A M \arrow[r] \arrow[r] \arrow[d] & (\goth{a}/I^n \goth{a}) \otimes_A M \cong (\goth{a} \otimes_A M) / I^n(\goth{a} \otimes_A M) \\
              M \arrow[r]                                         & M_k                                                                      &                         \end{tikzcd}. \]
          The $A$-module $\goth{a} / I^k \cap \goth{a}$ is a natural $A_k$-module, so we have the isomorphism
          \begin{equation*}
            (\goth{a} / I^k \cap \goth{a}) \otimes_A M \cong (\goth{a} / I^k \cap \goth{a}) \otimes_{A_k} M_k.
          \end{equation*}
          There is the natural isomorphism $\goth{a} / I^k \cap \goth{a} \cong (\goth{a} + I^k) / I^k$ (A.M., 2.1), which is an ideal in $A_k$.
          Since $M_k$ is a flat $A_k$-module, the map $(\goth{a} / I^k \cap \goth{a}) \otimes_A M \to M_k$ is injective.
          Hence, we get
          \begin{equation*}
            \ker{j} \subseteq \ker{f} \subseteq I^n(\goth{a} \otimes_A M).
          \end{equation*}

          The last algebraic result we need is the same statement of Lemma 3 but with the additional assumption that $A$ is an integral domain.
          In fact, a stronger statement holds, which we won't prove:
          \begin{lemma} [Nitsure, 5.11]
            Let $A$ be a Noetherian integeral domain, and $B$ an $A$-algebra of finite type.
            Then there exists $s \in A - \{ 0 \}$ such that $B_s$ is a free $A_s$-module.
          \end{lemma}

          We are ready to prove Lemma 3 by using the criterion of Lemma 4.
          The extended ideal $\goth{p}B_\goth{q}$ is contained in the maximal ideal of $B_{\goth{q}'}$ for any $\goth{q}' \in V(\goth{q})$, $B_{\goth{q}'}$ is flat over $A$ if and only if $B_{\goth{q}'} / \goth{p}B_{\goth{q}'}$ is flat over $A / \goth{p}$ and $\tor_1^A(B_{\goth{q}'}, A / \goth{p}) = 0$ (Lemma 5).
          Lemma 2 implies the isomorphism
          \begin{equation*}
            B_\goth{q} \otimes_B \tor_1^A(B, A / \goth{p}) \cong \tor_1^A(B_\goth{q}, A / \goth{p}) = 0,
          \end{equation*}
          and $\tor_1^A(B, A / \goth{p})$ is a finitely generated $B$-module because it is a $B$-submodule of $\goth{p} \otimes_A B$, so there exists $g \in B - \goth{q}$ such that
          \begin{equation*}
            B_g \otimes_B \tor_1^A(B, A / \goth{p}) \cong \tor_1^A(B_g, A / \goth{p}) = 0.
          \end{equation*}
          Thus, $\tor_1^A(B_{\goth{q}'}, A / \goth{p}) = 0$ for any $\goth{q}' \in V(\goth{q}) \cap X_g$.
          By Lemma 6, there exists $s \in A - \goth{p}$ such that $B_s / \goth{p}B_s$ is free over $A_s / \goth{p}A_s$, hence flat over $A / \goth{p}$ (Lemma 1).
          For any $\goth{q}' \in V(\goth{q}) \cap X_s$, $B_{\goth{q}'} / \goth{p}B_{\goth{q}'}$ is a localization of $B_{s} / \goth{p}B_s$, so it is flat over $A / \goth{p}$ as well.
          Hence, $V(\goth{q}) \cap X_g \cap X_s$ is a nonempty open subset of $V(\goth{q})$ contained in $E \cap V(\goth{q})$.
        \end{proof}

  \item[\textbf{11.}] Let $Y$ be a nonsingular curve of degree $d$ in $\PP_k^n$, over an algebraically closed field $k$. Show that
        \begin{equation*}
          0 \leq p_a(Y) \leq \frac{1}{2}(d - 1)(d - 2).
        \end{equation*}

        \begin{proof}
          If $n = 2$, the inequalitie becomes an equality even if $Y$ is singular or not irreducible (I, Ex. 7.2b) (Ex. 4.7), so assume $n > 2$.
          We claim there exists a linear subspace $L \subset \PP^n$ of dimension $n - 3$ such that $Y \cap L = \emptyset$.
          In fact, we can find a linear subspace of dimension $n - 2$ that does not meet $Y$, which is what we will prove.
          A general hyperplane $H \subset \PP^n$ intersects $Y$ at $d$ distinct points, and we want to find a hyperplane plane of $H$, which is a linear subspace of $\PP^n$ with dimension $n - 2$, that does not meet $Y \cap H$.
          Thus, the statement reduces to the following: for any set of distinct points $P_1, \dots, P_d \in \PP^n$, there exists a hyperplane $H \subset \PP^n$ such that $P_i \notin H$ for every $i$.
          Proceeding by induction on $n$, let $(x_0 : \cdots : x_n)$ be homogenous coordinates for $\PP^n$.
          The base case $n = 1$ is trivial because $k$ is infinite, so assume the statement for $\PP^{n - 1}$ with $n > 1$.
          Let $U_0$ be the affine patch of $\PP^n$ defined by $x_0 \neq 0$ so that $\PP^n$ is the disjoint union $U_0 \cup \PP^{n - 1}$.
          Relabel the points $P_1, \dots, P_{r}, Q_1, \dots, Q_s$ so that $P_i \in U_0$ and $Q_j \in \PP^{n - 1}$.
          Applying the inductive step to $\PP^{n - 1}$ and $\{Q_j\}$ gives a hyperplane $H$ defined by $$H : \quad f(x_1, \dots, x_n) = a_1x_1 + \cdots + a_nx_n = 0, \quad a_i \in k,$$ such that $Q_j \notin H$.
          Taking the projective cone of $H$ in $\PP^n$, we want to show there exists $a_0 \in k - \{0\}$ such that the hyperplane defined $a_0x_0 + f(x_1, \dots, x_n) = 0$ does not meet $P_i$ for every $i, \dots, r$.
          Indeed, if we consider $f(x_1, \dots, x_n)$ as a function on $U_0$, then we can pick any $a_0 \in k - \{0\}$ such that $a_0 \neq -f(P_i)$ for every $i = 1, \dots, r$.

          Now, after a suitable automorphism of $\PP^n$ we can assume $L$ corresponds to the linear subspace defined by $x_i = 0$ for $i = 0, 1, 2$.
          Consider the projection map $\pi : \PP^n - L \to \PP^2$ defined by $(x_0 : \cdots : x_n) \mapsto (x_0 : x_1 : x_2)$.
          For each $a \in k$, $a \neq 0$, consider the automorphism $\sigma_a$ of $\PP^n$ defined by
          \begin{align*}
            \sigma_a : \PP^n     & \to \PP^n,                                        \\
            (x_0 : \cdots : x_n) & \mapsto (x_0 : x_1 : x_2 : ax_3 : \cdots : ax_n).
          \end{align*}
          Further assume $\pi$ induces a birational from $Y$ to its image (I, Ex. 4.9).
          For each $a \neq 0$, let $Y_a = \sigma_a(Y)$ with $Y_1 = Y$.
          Then the $Y_a$ form a flat family parametrized by $\A^1 - \{0\}$, and according to (9.8) this family extends uniquely to a flat family defined over all of $\A^1$, and the fibre $Y_0$ over 0 agrees set-theoretically with the projection $\pi(Y) \subseteq \PP^2$.
          The degree and arithmetic genus of $Y_a$ is independent of $a$ (9.10), which means
          \begin{equation*}
            p_a(Y_0) = p_a(Y), \quad \deg{Y_0} = d.
          \end{equation*}
          If we take $Y_{(0)}$ to be $Y_0$ with its reduced induced structure, then we have an isomorphism of schemes $Y_{(0)} \cong \pi(Y)$.
          Since $\pi(Y)$ is a closed subscheme of $\PP^2$, one has
          \begin{equation*}
            p_a(Y_{(0)}) = \frac{1}{2}(d' - 1)(d' - 2),
          \end{equation*}
          where we take $d'$ to be the degree of $Y_{(0)}$ as a closed subscheme of $\PP^2$.
          The Hilbert polynomial of $Y_{(0)}$ as a closed subscheme of $\PP^2$ is same as $Y_{{(0)}}$ considered as a closed subcheme of $\PP^n$.
          This has the following implication: there is an exact sequence of $\fO_{Y_0}$-modules
          \[ \begin{tikzcd}
              0 \arrow[r] & \fI \arrow[r] & \fO_{Y_0} \arrow[r] & \fO_{Y_{(0)}} \arrow[r] & 0,
            \end{tikzcd} \]
          where $\fI$ is a nilpotent ideal sheaf of $\fO_{Y_0}$.
          Hilbert polynomials are additive (Ex. 5.1, 5.2a), so one has
          \begin{equation*}
            P_{Y_{0}}(n) = P_{\fI}(n) + P_{Y_{(0)}}(n),
          \end{equation*}
          where $P_{\fI}(n)$, $P_{Y_0}(n)$, and $P_{Y_{(0)}}(n)$ are the Hilbert polynomials of $\fI, \fO_{Y_0}$, and $\fO_{Y_{(0)}}$ with respect to $\fO_{\PP^n}(1)$, and by our remark, $d' = \deg P_{Y_{(0)}}(n)$.
          If we can show the following:
          \begin{itemize}
            \item[(1)] $p_a(Y_0) \leq p_a(Y_{(0)})$, and
            \item[(2)] $\deg {Y_{(0)}} \leq \deg{Y_0}$,
          \end{itemize}
          then we can combine the inequalities to conclude that
          \begin{equation*}
            0 \leq p_a(Y) = p_a(Y_0) \leq p_a(Y_{(0)}) = \frac{1}{2}(d' - 1)(d' - 2) \leq \frac{1}{2}(d - 1)(d - 2).
          \end{equation*}
          The degree of the Hilbert polynomial is at most the dimension the dimension of the support, and its leading coefficient is always positive, which proves (2).
          In addition, we can plug in $n = 0$ to obtain the equality
          \begin{equation*}
            P_\fI(0) + p_a(Y_0) = p_a(Y_{(0)}),
          \end{equation*}
          so to prove (1), we want to show $P_\fI(0) = \chi(\fI) \geq 0$.
          \textit{Incomplete.}
        \end{proof}

        \newpage

        \begin{lemma} [Stacks Project, Tag 00HT]
          Let $A \to B$ be a ring homomorphism, and let $M$ be a $B$-module.
          The following are equivalent:
          \begin{itemize}
            \item[(i)] $M$ is a flat $A$-module;
            \item[(ii)] for every prime ideal $\goth{q}$ of $B$, $\goth{p}$ its contraction in $A$, $M_\goth{q}$ is a flat $A_\goth{p}$-module.
          \end{itemize}
        \end{lemma}

        \begin{proof}
          $(i) \implies (ii)$ Suppose $M$ is flat as an $A$-module.
          Let $\goth{q}$ be a prime ideal of $B$, and let $\goth{p}$ be its contraction in $A$.
          Then $\goth{a} \otimes_A M \to M$ is injective for any ideal $\goth{a}$ of $A$.
          We want to show $\goth{a}A_\goth{p} \otimes_{A_\goth{p}} M_\goth{q} \to M_\goth{q}$ is an $A_\goth{p}$-module monomorphism.
          By (A.M., 2.15), and because $M_\goth{q}$ is a natural $A_\goth{p}$-module, we have the isomorphisms
          \begin{align*}
            (\goth{a} \otimes_A M) \otimes_B B_\goth{q} & \cong \goth{a} \otimes_A (M \otimes_B B_\goth{q})                     \\
                                                        & \cong \goth{a} \otimes_A (A_\goth{p} \otimes_{A_\goth{p}} M_\goth{q}) \\
                                                        & \cong \goth{a}A_\goth{p} \otimes_{A_\goth{p}} M_\goth{q},
          \end{align*}
          which shows $\goth{a} A_\goth{p} \otimes_{A_\goth{p}} M_\goth{q} \to M_\goth{q}$ is given by localizing the $B$-module homomorphism $\goth{a} \otimes_A M \to M$ at the prime ideal $\goth{q}$.
          Localizing preserves injections, so we are done.

          \noindent $(ii) \implies (i)$ We want to show $\goth{a} \otimes_A M \to M$ is injective for any ideal $\goth{a}$ of $A$.
          Because $M$ is a $B$-module, localizing at $\goth{q}$ gives $(\goth{a} \otimes_A M) \otimes_B B_\goth{q} \to M_\goth{q}$, which is injective by assumption and isomorphism of above.
          Thus, $\goth{a} \otimes_A M \to M$ is injective as a $B$-module homomorphism, and therefore injective as an $A$-module homomorphism.
        \end{proof}

        \begin{lemma}
          Let $A \to B \to C$ be ring homomorphisms, and suppose $C$ is flat over $B$.
          Then for every ideal $\goth{a} \subseteq A$, there is a natural isomorphism
          \begin{equation*}
            C \otimes_B \tor_1^A(B, A / \goth{a}) \cong \tor_1^A(C, A / \goth{a}).
          \end{equation*}
        \end{lemma}

        \begin{proof}
          Consider the exact sequence
          \[ \begin{tikzcd}
              0 \arrow[r] & \goth{a} \arrow[r] & A \arrow[r] & A/\goth{a} \arrow[r] & 0.
            \end{tikzcd} \]
          Tensoring by $B$ over $A$ gives an exact sequence with $\tor$
          \[ \begin{tikzcd}
              0\arrow[r] & \tor_1^A(B, A/\goth{a}) \arrow[r] & \goth{a} \otimes_A B \arrow[r] & B \arrow[r] & B/\goth{a}B \arrow[r] & 0.
            \end{tikzcd} \]
          Each term in the sequence is a natural $B$-module, so we can tensor the sequence by $C$ over $B$
          \[ \begin{tikzcd}
              0\arrow[r] & C \otimes_B \tor_1^A(B, A/\goth{a}) \arrow[r] & \goth{a} \otimes_A C \arrow[r] & C \arrow[r] & C/\goth{a}C \arrow[r] & 0.
            \end{tikzcd} \]
          Here, we used the fact that
          \begin{equation*}
            (\goth{a} \otimes_A B) \otimes_B C \cong \goth{a} \otimes_A (B \otimes_B C) \cong \goth{a} \otimes_A C.
          \end{equation*}
          We can also tensor the original sequence by $C$ over $A$ to get
          \[ \begin{tikzcd}
              0\arrow[r] & \tor_1^A(C, A/\goth{a}) \arrow[r] & \goth{a} \otimes_A C \arrow[r] & C \arrow[r] & C/\goth{a}C \arrow[r] & 0.
            \end{tikzcd} \]
          The $A$-module homomorphism $B \to C$ induces $\tor_1^A(B, A / \goth{a}) \to \tor_1^A(C, A / \goth{a})$, and $\tor_1^A(C, A)$ is a natural $C$-module, so there is a natural map
          \begin{equation*}
            \phi : C \otimes_B \tor_1^A(B, A / \goth{a}) \to \tor_1^A(C, A / \goth{a})
          \end{equation*}
          given by multiplication.
          Consider the commutative diagram
          \[ \begin{tikzcd}
              0 \arrow[r] & {C \otimes_B \tor_1^A(B, A / \goth{a})} \arrow[r] \arrow[d, "\phi"] & \goth{a} \otimes_A C \arrow[r] \arrow[d, Rightarrow, no head] & C \arrow[r] \arrow[d, Rightarrow, no head] & C/\goth{a}C \arrow[r] \arrow[d, Rightarrow, no head] & 0 \\
              0 \arrow[r] & {\tor_1^A(C, A / \goth{a})} \arrow[r]                       & \goth{a} \otimes_A C \arrow[r]                                & C \arrow[r]                                & C / \goth{a}C \arrow[r]                              & 0
            \end{tikzcd}. \]
          The five lemma implies $\phi$ is actually an isomorphism.
        \end{proof}

\end{enumerate}

\end{document}

% \item A flat morphism $f : X \to Y$ of finite type of Noetherian schemes is open, i.e., for every open subset $U \subseteq X$, $f(U)$ is open in $Y$.

% \item Do the calculation of (9.8.4) for the curves of (I, Ex. 3.14). Show that you get an embedded point at the cusp of the plane cubic curve.

% \item \textit{Very Flat Families.} For any closed subscheme $X \subseteq \PP^n$, we denote by $C(X) \subseteq \PP^{n + 1}$ the projective cone over $X$ (I, Ex. 2.10). If $I = \subseteq k[x_0, \dots, x_n]$ is the (largest) homogenous ideal of $X$, then $C(X)$ is defined by the ideal generated by $I$ in $k[x_0, \dots, x_{n + 1}]$.
% \begin{itemize}
%   \item[(a)] Give an example to show that if $\{ X_t \}$ is a flat family of closed subschemes of $\PP^n$, then $\{ C(X_t) \}$ need not be a flat family in $\PP^{n + 1}$.
%   \item[(b)] To remedy this situation, we make the following definition. Let $X \subseteq \PP_T^n$ be a closed subscheme, where $T$ is a Noetherian integral scheme. For each $t \in T$, let $I_t \subseteq S_t = k(t)[x_0, \dots, x_n]$ be the homogenous ideal of $X_t$ in $\PP_{k(t)}^n$. We say that the family $\{X_t \}$ is \textit{very flat} if for all $d \geq 0$,
%   \begin{equation*}
%     \dim_{k(t)}(S_t / I_t)_d
%   \end{equation*}
%   is independent of $t$.
%   \item[(c)] If $\{X_t\}$ is a very flat family in $\PP^n$, show that it is flat. Show also that $\{ C(X_t) \}$ is a very flat family in $\PP^{n + 1}$, and hence flat.
%   \item[(d)] If $\{ X_{(t)}\}$ is an algebraic family of projectively normal varieties in $\PP_k^n$, parametrized by a nonsingular curve $T$ over an algebraically closed field $k$, then $\{X_{(t)} \}$ is a very flat family of schemes.
% \end{itemize}

% \item Let $Y \subseteq \PP^n$ be a nonsingular variety of dimension $\geq 2$ over an algebraically closed field $k$. Suppose $\PP^{n - 1}$ is a hyperplane in $\PP^n$ which does not contain $Y$, and such that the scheme $Y' = Y \cap \PP^{n - 1}$ is also nonsingular. Prove that $Y$ is a complete intersection in $\PP^n$ if and only if $Y'$ is a complete intersection in $\PP^{n - 1}$.

% \item Let $Y \subseteq X$ be a closed subscheme, where $X$ is a scheme of finite type over a field $k$. Let $D = k[t] / (t^2)$ be the ring of dual numbers, and define an \textit{infinitesimal deformation} of $Y$ \textit{as a closed subscheme of $X$}, to be a closed subscheme $Y' \subseteq X \times_k D$, which is flat over $D$, and whose closed fiber is $Y$. Show that these $Y'$ are classified by $H^0(Y, \fN_{Y/X})$, where
% \begin{equation*}
%   \fN_{Y/X} = \rHom_{\fO_Y}(\mathscr{J}_y/\mathscr{J}_Y^2, \fO_Y).
% \end{equation*}



% \item A $k$-algebra $A$ is said to be \textit{rigid} if it has no infinitesimal deformations, or equivalently, by (Ex. 9.8) if $T^1(A) = 0$. Let $A = k[x, y, z, w] / (x, y) \cap (z, w)$, and show that $A$ is rigid. This corresponds to two planes in $\A^4$ which meet at a point.

% \item A scheme $X_0$ over a field $k$ is \textit{rigid} if it has no infinitesimal deformations.
% \begin{itemize}
%   \item[(a)] Show that $\PP_k^1$ is rigid, using (9.13.2).
%   \item[(b)] one might think that if $X_0$ is rigid over $k$, then every global deformation of $X_0$ is locally trivial. Show that this is not so, by constructing a proper, flat morphism $f : X \to \A^2$ over $k$ algebraically closed, such that $X_0 \cong \PP_k^1$, but there is no open neighborhood $U$ of $0$ in $\A^2$ for which $f^{-1}(U) \cong U \times \PP^1$.
%   \item[(c)] Show, however, that one can trivialize a global deformation of $\PP^1$ after a flat base extension, in the following sense: let $f : X \to T$ be a flat projective morphism, where $T$ is a nonsingular curve over $k$ algebraically closed. Assume there is a closed point $t \in T$ such that $X_t \cong \PP_k^1$. Then there exists a nonsingular curve $T'$, and a flat morphism $g : T' \to T$, whose image contains $t$, such that if $X' = X \times_T T'$ is the base extension, then the new family $f' : X' \to T'$ is isomorphic to $\PP_T^1 \to T'$.
% \end{itemize}

% \item[\textbf{8.}] Let $A$ be a finitely generated $k$-algebra. Write $A$ as a quotient of a polynomial ring $P$ over $k$, and let $J$ be the kernel:
% \begin{equation*}
% 0 \to J \to P \to A \to 0.
% \end{equation*}
% Consider the exact sequence of (II, 8.4A)
% \begin{equation*}
% J / J^2 \to \Omega_{P/k} \otimes_P A \to \Omega_{A / k} \to 0.
% \end{equation*}
% Apply the functor $\text{Hom}_A(\cdot, A)$ and let $T^1(A)$ be the cokernel:
% \begin{equation*}
% \Hom_A(\Omega_{P/k} \otimes A, A) \to \Hom_A(J/J^2, A) \to T^1(A) \to 0.
% \end{equation*}
% Now use the construction of (II, Ex. 8.6) to show that $T^1(A)$ classifies infinitesimal deformations of $A$, i.e., algebras $A'$ flat over $D = k[t] / (t^2)$, with $A' \otimes_D k \cong A$. It follows that $T^1(A)$ is independent of the given representation of $A$ as quotient of a polynomial ring $P$.